\documentclass[11pt]{article}
\usepackage[fleqn]{amsmath}
\usepackage{amssymb, amsthm}
\usepackage[utf8]{inputenc}
\usepackage[margin=1in]{geometry}
\usepackage{graphicx}
\usepackage{enumitem}
\theoremstyle{plain}
\newtheorem{proposition}{Proposition}
\newcommand{\fit}[1]{\resizebox{\ifdim\width>\linewidth\linewidth\else\width\fi}{!}{#1}}
\setlength{\parindent}{0pt}
\begin{document}
\begin{proposition}[smith-diagonalization]
Suppose $a \text{ s.t. } \operatorname{is\text{-}euclidean\text{-}ring}(a)$. Let $k \in \mathbb{N}, n \in \mathbb{N}, p \in \operatorname{mat}(k, n, \operatorname{carr}(a))$. There is $d$ such that 
\begin{equation*}
\begin{array}{@{}l@{}} \operatorname{mat\text{-}equiv}(a, k, n, p, d) \wedge \\ \operatorname{is\text{-}diagonal}(a, k, n, d) \end{array}
\end{equation*}

\end{proposition}

\begin{proof}
\begin{itemize}[leftmargin=2.8em, itemsep=4pt, parsep=1pt]
  \item[\textbf{3}] By $\operatorname{euclidean\text{-}ring\text{-}is\text{-}integral\text{-}domain}$:
\begin{equation*}
\operatorname{is\text{-}integral\text{-}domain}(a)
\end{equation*}

  \item[\textbf{4}] By $\operatorname{integral\text{-}domain\text{-}is\text{-}commutative\text{-}ring}$:
\begin{equation*}
\operatorname{is\text{-}commutative\text{-}ring}(a)
\end{equation*}

  \item[\textbf{5}] By $\operatorname{commutative\text{-}ring\text{-}is\text{-}ring}$:
\begin{equation*}
\operatorname{is\text{-}ring}(a)
\end{equation*}

  \item[\textbf{6}] By induction on $k$.
  \item[]\textbf{Base case} ($k = 0$):
  \item[\textbf{7--10}] Take $w :=$ $p$.
  \item[\textbf{11--13}] By $\operatorname{mat\text{-}equiv\text{-}refl}$:
\begin{equation*}
\operatorname{mat\text{-}equiv}(a, 0, n, p, p)
\end{equation*}
  \emph{(holds by assumption)}
  \item[\textbf{14}] By $\operatorname{is\text{-}diagonal}$, reduce to
\begin{equation*}
\begin{array}{@{}l@{}} \forall i. \forall j. \\ \quad (i \in \operatorname{interval}(1, 0)) \Rightarrow  \\ \quad \quad (j \in \operatorname{interval}(1, n)) \Rightarrow  \\ \quad \quad \quad \neg (i = j) \Rightarrow  \\ \quad \quad \quad \quad (\operatorname{entry}(p, i, j) = \operatorname{zero}(a)) \end{array}
\end{equation*}

  \item[\textbf{15--19}] By $\operatorname{interval\text{-}lo}$:
\begin{equation*}
(1 \leq i)
\end{equation*}

  \item[\textbf{20}] By $\operatorname{interval\text{-}hi}$:
\begin{equation*}
(i \leq 0)
\end{equation*}

  \item[\textbf{21}] \emph{Establish in-carrier / typing side-conditions:}
\begin{equation*}
\begin{array}{@{}l@{}}
(i \in \mathbb{N})
\end{array}
\end{equation*}

  \item[\textbf{22}] By $\operatorname{nn\text{-}not\text{-}le\text{-}zero\text{-}pos}$:
\begin{equation*}
\neg (i \leq 0)
\end{equation*}

  \item[]\textbf{Inductive step} ($k \to k + 1$):
  \item[\textbf{23--41}] \emph{Split into cases} on whether
\begin{equation*}
(1 \leq n)
\end{equation*}

  \item[]\emph{Case} $\lnot (1 \leq n)$:
  \item[\textbf{42}] Take $d :=$ $p$.
  \item[\textbf{43--45}] By $\operatorname{mat\text{-}equiv\text{-}refl}$:
\begin{equation*}
\operatorname{mat\text{-}equiv}(a, (k + 1), n, p, p)
\end{equation*}
  \emph{(holds by assumption)}
  \item[\textbf{46}] By $\operatorname{is\text{-}diagonal}$, reduce to
\begin{equation*}
\begin{array}{@{}l@{}} \forall i. \forall j. \\ \quad (i \in \operatorname{interval}(1, (k + 1))) \Rightarrow  \\ \quad \quad (j \in \operatorname{interval}(1, n)) \Rightarrow  \\ \quad \quad \quad \neg (i = j) \Rightarrow  \\ \quad \quad \quad \quad (\operatorname{entry}(p, i, j) = \operatorname{zero}(a)) \end{array}
\end{equation*}

  \item[\textbf{47--51}] By $\operatorname{interval\text{-}lo}$:
\begin{equation*}
(1 \leq j)
\end{equation*}

  \item[\textbf{52}] By $\operatorname{interval\text{-}hi}$:
\begin{equation*}
(j \leq n)
\end{equation*}

  \item[\textbf{53}] By $\operatorname{nn\text{-}le\text{-}trans}$:
\begin{equation*}
(1 \leq n)
\end{equation*}

  \item[]\emph{Case} $(1 \leq n)$:
  \item[\textbf{54--55}] By $\operatorname{nn\text{-}pos\text{-}is\text{-}succ}$:
\begin{equation*}
\begin{array}{@{}l@{}} \exists q \in \mathbb{N}. \\ \quad (n = (q + 1)) \end{array}
\end{equation*}

  \item[\textbf{56--58}] By $\operatorname{eq\text{-}sym}$:
\begin{equation*}
((q + 1) = n)
\end{equation*}

  \item[\textbf{59}] Substitute the established equation.
  \item[\textbf{60}] By a cut on the auxiliary claim.
  \item[\textbf{61--62}] Substitute the established equation.  \emph{(holds by assumption)}
  \item[\textbf{63--75}] \emph{Split into cases} on whether
\begin{equation*}
\begin{array}{@{}l@{}} \exists \mathit{i0}. \exists \mathit{j0}. \\ \quad (\mathit{i0} \in \operatorname{interval}(1, (k + 1))) \wedge \\ \quad (\mathit{j0} \in \operatorname{interval}(1, (q + 1))) \wedge \\ \quad \neg (\operatorname{entry}(p, \mathit{i0}, \mathit{j0}) = \operatorname{zero}(a)) \end{array}
\end{equation*}

  \item[]\emph{Case} $\exists \mathit{i0}.\; \exists \mathit{j0}.\; ((\mathit{i0} \in \operatorname{interval}(1, (k + 1))) \wedge ((\mathit{j0} \in \operatorname{interval}(1, (q + 1))) \wedge \neg (\operatorname{entry}(p, \mathit{i0}, \mathit{j0}) = \operatorname{zero}(a))))$:
  \item[\textbf{76}] By $\operatorname{clear\text{-}pivot\text{-}cross}$:
\begin{equation*}
\begin{array}{@{}l@{}} \exists \mathit{c2}. \\ \quad \operatorname{mat\text{-}equiv}(a, (k + 1), (q + 1), p, \mathit{c2}) \wedge \\ \quad \mathit{c2} \\ \quad =\; \operatorname{border}(a, \operatorname{entry}(\mathit{c2}, 1, 1), \operatorname{submat}(\mathit{c2}, k, q), k, q) \wedge \\ \quad \neg (\operatorname{entry}(\mathit{c2}, 1, 1) = \operatorname{zero}(a)) \end{array}
\end{equation*}

  \item[\textbf{77--81}] \emph{Establish in-carrier / typing side-conditions:}
\begin{equation*}
\begin{array}{@{}l@{}}
(\mathit{c2} \in \operatorname{mat}((k + 1), (q + 1), \operatorname{carr}(a))) \\
(\operatorname{submat}(\mathit{c2}, k, q) \in \operatorname{mat}(k, q, \operatorname{carr}(a)))
\end{array}
\end{equation*}

  \item[\textbf{82--83}] Instantiate the hypothesis:
\begin{equation*}
\begin{array}{@{}l@{}} \exists d. \\ \quad \operatorname{mat\text{-}equiv}(a, k, q, \operatorname{submat}(\mathit{c2}, k, q), d) \wedge \\ \quad \operatorname{is\text{-}diagonal}(a, k, q, d) \end{array}
\end{equation*}

  \item[\textbf{84--88}] \emph{Establish in-carrier / typing side-conditions:}
\begin{equation*}
\begin{array}{@{}l@{}}
(1 \in \operatorname{interval}(1, (k + 1))) \\
(1 \in \operatorname{interval}(1, (q + 1))) \\
(\operatorname{entry}(\mathit{c2}, 1, 1) \in \operatorname{carr}(a))
\end{array}
\end{equation*}

  \item[\textbf{89}] By $\operatorname{bordering}$:
\begin{equation*}
\operatorname{mat\text{-}equiv}(a, (k + 1), (q + 1), \mathit{c2}, \operatorname{border}(a, \operatorname{entry}(\mathit{c2}, 1, 1), d, k, q))
\end{equation*}

  \item[\textbf{90}] By $\operatorname{border\text{-}is\text{-}diagonal}$:
\begin{equation*}
\operatorname{is\text{-}diagonal}(a, (k + 1), (q + 1), \operatorname{border}(a, \operatorname{entry}(\mathit{c2}, 1, 1), d, k, q))
\end{equation*}

  \item[\textbf{91}] By $\operatorname{mat\text{-}equiv\text{-}trans}$:
\begin{equation*}
\operatorname{mat\text{-}equiv}(a, (k + 1), (q + 1), p, \operatorname{border}(a, \operatorname{entry}(\mathit{c2}, 1, 1), d, k, q))
\end{equation*}

  \item[\textbf{92--95}] Take $d :=$
\begin{equation*}
\operatorname{border}(a, \operatorname{entry}(\mathit{c2}, 1, 1), d, k, q)
\end{equation*}
  \emph{(holds by assumption)}
  \item[]\emph{Case} $\lnot \exists \mathit{i0}.\; \exists \mathit{j0}.\; ((\mathit{i0} \in \operatorname{interval}(1, (k + 1))) \wedge ((\mathit{j0} \in \operatorname{interval}(1, (q + 1))) \wedge \neg (\operatorname{entry}(p, \mathit{i0}, \mathit{j0}) = \operatorname{zero}(a))))$:
  \item[\textbf{96}] Take $d :=$ $p$.
  \item[\textbf{97--99}] By $\operatorname{mat\text{-}equiv\text{-}refl}$:
\begin{equation*}
\operatorname{mat\text{-}equiv}(a, (k + 1), (q + 1), p, p)
\end{equation*}
  \emph{(holds by assumption)}
  \item[\textbf{100}] By $\operatorname{is\text{-}diagonal}$, reduce to
\begin{equation*}
\begin{array}{@{}l@{}} \forall i. \forall j. \\ \quad (i \in \operatorname{interval}(1, (k + 1))) \Rightarrow  \\ \quad \quad (j \in \operatorname{interval}(1, (q + 1))) \Rightarrow  \\ \quad \quad \quad \neg (i = j) \Rightarrow  \\ \quad \quad \quad \quad (\operatorname{entry}(p, i, j) = \operatorname{zero}(a)) \end{array}
\end{equation*}

  \item[\textbf{101--118}] \emph{Split into cases} on whether
\begin{equation*}
(\operatorname{entry}(p, i, j) = \operatorname{zero}(a))
\end{equation*}
  \emph{(holds by assumption)}
  \item[]\emph{Case} $\lnot (\operatorname{entry}(p, i, j) = \operatorname{zero}(a))$:
  \item[\textbf{119}] By a cut on the auxiliary claim.
  \item[\textbf{120}] Take $\mathit{i0} :=$ $i$.
  \item[\textbf{121--126}] Take $\mathit{j0} :=$ $j$.  \emph{(holds by assumption)}
  \item[\textbf{127}] 
\end{itemize}
\end{proof}
\end{document}
